\documentclass{rapport}
\usepackage{lipsum}
\usepackage{dirtytalk}
\usepackage[style=ieee]{biblatex} %Imports biblatex package
\usepackage{csquotes}
\usepackage{amsmath}
\usepackage{tcolorbox}
\usepackage{hyperref}
\usepackage{pdfpages}
\usepackage{multirow}

\hypersetup{
    colorlinks=true,
    linkcolor=blue,
    filecolor=magenta,      
    urlcolor=cyan,
    }
\urlstyle{same}

\newlength\tindent
\setlength{\tindent}{\parindent}
\setlength{\parindent}{0pt}
\renewcommand{\indent}{\hspace*{\tindent}}

% custom commands
\newcommand{\newparagraph}{\vspace{0.5em}}
\newcommand{\itembullet}{\item[$\bullet$]}

\addbibresource{sources.bib} %Import the bibliography file
\title{LSA_Lab_01_Template_Rapport_Première_Manip} %Titre du fichier

\begin{document}

%----------- Informations du rapport ---------

\logo{img/henallux.jpg}
\dep{IESN}
\titre{Lab \#01 - Template} %Titre du fichier .pdf
\cours{NLST1} %Nom du cours
\sujet{Laboratoire de sciences appliquées à l'informatique} %Nom du sujet

\enseignant{Prénom \textsc{nom}} %Nom de l'enseignant

\eleves{Prénom \textsc{Nom} \\ %Nom étudiants
        Prénom \textsc{Nom} \\
        Prénom \textsc{Nom} \\
        Prénom \textsc{Nom} \\
        Prénom \textsc{Nom} }
       

%----------- Initialisation -------------------
        
\fairemarges %Afficher les marges
\fairepagedegarde %Créer la page de garde
\tabledematieres %Créer la table de matières

%------------ Corps du rapport ----------------

\include{contenu/1_Introduction}

\newpage

\include{contenu/2_Rappels_théoriques}

\newpage

\include{contenu/3_Manipulation_pratique}

\newpage

\include{contenu/4_Conclusion}

%------------ Bibliographie ----------------

\printbibliography

% %------------- Commandes utiles ----------------

% \section{Quelques commandes}

% Voici quelques commandes utiles :

% %------ Pour insérer et citer une image centralisée -----

% \insererfigure{img/henallux.jpg}{3cm}{Légende de la figure}{Label de la figure}
% % Le premier argument est le chemin pour la photo
% % Le deuxième est la hauteur de la photo
% % Le troisième la légende
% % Le quatrième le label
% Ici, je cite l'image \ref{fig: Label de la figure}

% %------- Pour insérer et citer une équation --------------

% \begin{equation} \label{eq: exemple}
% \rho + \Delta = 42
% \end{equation}

% L'équation \ref{eq: exemple} est cité ici. 

% % ------- Pour écrire des variables ----------------------

% Pour écrire des variables dans le texte, il suffit de mettre le symbole \$ entre le texte souhaité comme : constante $\rho$. 


\end{document}
